\section{Introduction}\label{c:introduction}
A boundary-to-boundary walk can be weighted critically without fixing its
length. In that law the separation of the endpoints supplies the spatial
scale. Our question is whether prescribing the terminal boundary point,
rather than summing over an interval, changes the power relating length
to that scale. The distinction is substantial: a macroscopic terminal
interval contains order $R$ ports, and a bound on its total exceptional
mass does not bound the exceptional mass at one specified port.

We use the regular honeycomb lattice, embedded as the centers of the
triangles of an equilateral triangular tiling of side one. Neighboring
centers are joined across a common triangle edge. A \emph{port} is the
midpoint of a boundary edge of a domain made of whole triangles. A path
enters and leaves normally at its ports, visits only centers in the
interior, and visits no center twice. Its length $L(\gamma)$ is the number
of visited centers, and its critical weight is
\[
 w(\gamma)=\kappa^{L(\gamma)},\qquad
 \kappa=(2\cos(\pi/8))^{-1}=(2+\sqrt2)^{-1/2}.
\]
Ports carry no weight. For a domain $D$ and boundary ports $a,b$, set
\[
 K_D(a,b)=\sum_{\gamma:a\to b\text{ in }D}w(\gamma).
\]
Whenever the sum is positive and finite, dividing by it defines the
critical chord law. The path has no other boundary exit. A \emph{pure}
cut follows one of the three tiling-edge directions; its admissible
translates bound whole triangular rows. Distances refer to this fixed
embedding.

\begin{theorem}[Prescribed boundary endpoints]\label{c:main}
Let $D_R$ be convex pure lattice hexagons which, after scaling by $R$,
converge to a regular hexagon. Let $a_R,b_R$ be boundary ports whose
scaled limits are distinct points in the relative interiors of sides.
Then, for every $\varepsilon>0$, the critical chord law satisfies
\[
 \mathbb P_{D_R;a_R,b_R}
 \{R^{4/3-\varepsilon}\le L(\gamma)\le R^{4/3+\varepsilon}\}
 \longrightarrow1.
\]
The same conclusion holds in a full half-plane for prescribed boundary
ports at distance $R$, in a strip of height comparable to $R$ for
one prescribed port on each of the two opposite boundary cuts, at mutual
distance $O(R)$, and in a wedge of angle
$\pi/3$ or $2\pi/3$ formed by two nonparallel pure cuts, with one prescribed port on each ray.
In the wedge, each endpoint has clearance comparable to $R$ from the other cut and
distance $O(R)$ from the apex. Constants are uniform on fixed compact
sets of these nondegenerate scaled geometries. In each case the
normalizing mass is at least $cR^{-5/4}$.
\end{theorem}

The theorem concerns each admissible choice of the two ports. A first
step, proved in Section~\ref{c:window}, fixes the source and sums the
terminal port over a separated macroscopic interval. Its normalizing
mass is of order $R^{-1/4}$. This intermediate law is also useful in its
own right, and remains distinct from the prescribed-endpoint law.

For half-plane arches we additionally prove
$\mathbb E L=R^{4/3+o(1)}$. This requires control of the weighted upper
tail; convergence in probability alone would not imply it. We also record
the terminal-summed mean law from the proof of \cite[Theorem 8.1]{Renewal},
where the gap ranges over $[R,2R]$ and the arch diameter is confined to a
sufficiently large fixed multiple of $R$. For strip
bridges, a second argument proves a mean law uniformly over sublinear
terminal offsets on a single set of heights of natural density one.
It also proves the same mean exponent in the critical mixture of aligned
bridges at all even heights in $[H,2H]$. These mean statements retain their
own height and weighting conventions.

\paragraph{Historical context.}
Nienhuis's analysis of the dilute $O(n)$ model predicted the planar
self-avoiding-walk exponents, including the length--distance power
$4/3$~\cite{Nienhuis1982}. Duminil-Copin and Smirnov proved the honeycomb
connective constant using a local winding relation and a boundary
identity~\cite{DuminilCopinSmirnov2012}. Their mid-edge convention is the
port convention used here. For critical strip mass, Beaton, Bousquet-M\'elou, de Gier,
Duminil-Copin and Guttmann proved decay to zero in their work on surface
adsorption~\cite{BeatonEtAl2014}. Glazman and Manolescu gave a shorter
proof and a logarithmic subsequence bound, and proved boundary two-point
invariance for columnwise rhombic half-planes with angles in
$[\pi/3,2\pi/3]$~\cite{GlazmanManolescu2017}. Krachun and Panagiotis
subsequently proved a polynomial upper bound for strip-crossing mass
and a quantitative sub-ballistic estimate for uniform honeycomb
walks~\cite{KrachunPanagiotis2026}. These results distinguish critical
activity, decay of crossing mass, and length--distance exponents;
the finite inputs below specify the stronger estimates used here.

The irreducible-bridge construction goes back to
Kesten~\cite{Kesten1963}; its honeycomb mid-edge form appears in
\cite[Appendix]{BeatonEtAl2014}. On the square lattice, the infinite
critical law and its relation to strip-spanning laws are developed by
Lawler--Schramm--Werner~\cite[Appendix]{LawlerSchrammWerner2004} and
Dyhr--Gilbert--Kennedy--Lawler--Passon~\cite[Section 2.4]{DyhrEtAl2011}.
Gilbert's appendix records Kennedy's heuristic connecting strip mass
of order $h^{-1/4}$ to an irreducible height tail of order $h^{-3/4}$
and renewal heights of order $n^{4/3}$~\cite[Appendix]{Gilbert2015}.
Here whole irreducible bridges are independent increments; height,
length and lateral displacement within one increment remain dependent.

The quantitative starting point is the sharp finite marked-polygon and
bridge estimates of \cite{UniformMarked}, together with the renewal
estimates of \cite{Renewal}; Section~\ref{c:inputs} states exactly what
we use. These are stronger inputs than the preceding critical and
sub-ballistic results. The terminal-free fixed-height probability law
and the prescribed half-plane arch probability law already appear in
\cite[Propositions 5.7 and 5.8]{Renewal}. We give a different proof of the
latter, through common-level pair renewals and arch-kernel regularity,
and develop the local estimates needed to assemble the prescribed-endpoint
laws in the other geometries. The mean estimates require additional
length-weighted control and retain separate finite inputs.

\paragraph{The three endpoint mechanisms.}
For two ports on one line, grow independent bridge sequences upward from
the two ports and stop at their common renewal levels. Their common
height transform has exponent $1/2$. A separate lateral window estimate
gives $O(m^{-2})$ atoms for the difference after $m$ common blocks. This supplies both regularity
of the arch kernel and control of a long final joining arch; the complete
argument is in Section~\ref{c:same-side}.

For parallel cuts one must match height and lateral displacement at once.
A height-tail bound alone does not pay for one prescribed height.
Section~\ref{c:smoothness-section} bounds all fixed-order differences of
the strip mass, starting from the finite projection representation in
\cite{UniformMarked}. This local regularity gives the one-piece height
atom bound in Section~\ref{c:atoms-section}. Together with two-dimensional
concentration and a sum of heightwise maximal atoms, it pays the parallel
endpoint restriction in Section~\ref{c:parallel-section}.
For nonparallel cuts, small normal prefixes
initially exhaust only a sum over terminal ports. A local visit estimate
and an artificial-arc boundary identity then control the exceptional
mass at the prescribed terminal port. Each case thus pays for its actual
endpoint restriction before dividing by the chord mass;
Section~\ref{c:tilted} completes the nonparallel case.

The means use two further arguments. In Section~\ref{c:mean}, the
calibrated finite estimates of \cite{Calibrated} and their renewal
consequences control marked arms at both ends of an arch. The resulting
length-weighted contributions are summable over arch heights at each
prescribed gap. In Section~\ref{c:pvmeans},
the finite strip moments of \cite{PolynomialVacuum} feed a renewal-reward
identity. A two-bridge boundary comparison controls variation between
terminal ports after averaging over heights, leading to the density-one
and aligned-mixture mean laws rather than an all-height pointwise claim.
